%
% $Log: instguide.tex,v $
% Revision 1.6  2002/10/21 15:26:44  mjm
% to ohio, revised
%
% Revision 1.5  2001/08/15 18:07:03  mjm
% APPM preprint 466
%
% Revision 1.4  1999/08/09 19:55:27  mjm
% beta release
%
% Revision 1.3  1999/04/06 19:29:00  mjm
% Matts comments
%
% Revision 1.2  1999/04/06 19:22:02  mjm
% Ellens comments
%
%

%%%%%%%%%%%%%%%%%%%%%%%%%%%%%%%%%%%%%%%%%%%%%%%%%%%%%%%%%%%%%%%%%%%%%%%%%
\handout{Instructor's Guide to Good Problems}

 This guide is for the main instructor, meaning the person in charge of
organizing the course. It lists those tasks that the instructor must
perform to implement this scheme, indicates when they must be performed,
and estimates how much time they are expected to take.
If there is more than one instructor, then several of these tasks can be
shared.

The first time you use this method, there are extra tasks, which
we list first.
\begin{enumerate}
	\item Read this packet, learn this material, make modifications,
and decide that it is a good thing to do. This should be performed well
before the term starts, to allow the ideas to settle. Time required: 3
hours.
	\item Convince the other people involved with the course to go
along with this plan. This task should be completed a month before the
term starts. Time required: 1 hour.
	\item Deal with unanticipated problems. Time required: 2 hours.
\end{enumerate}

Each time the method is used, there are the following tasks.
\begin{enumerate}
\item 
Choose the good problems, schedule them, and decide on the time
schedule for each new skill. The problems should not be harder than
the regular homework problems, but should be substantive enough that
there is something to write about. They must also be chosen to be
relevant to the new skill learned that week. (If the handout is on
graphing, the problem must require a graph.) Since there are only six
handouts, there is flexibility in
scheduling certain skills for certain sections, or avoiding busy
(test) weeks.  You might consider making the good problem on the
previous week's material, since a student will not be able to write
about new material that they do not yet understand.  This task should
be performed when the regular homework problems are chosen. Time
required: 1 hour.  \item Provide initial training for teaching
assistants. You must introduce this method, distribute materials and
discuss its implementation with the teaching assistants. This should
occur one week before the term starts. Time required: 2 hours.
\item Explain the method to the students. The concept should be
introduced and the student's guide distributed in the first week of
classes. The skills handouts could be distributed all at once at the
beginning or as they are used. Time required: 10 minutes of class
time.  \item Provide feedback on the handouts. Any deficiencies in the
method or the handouts should be noted. At the end of the term
comments from the instructors, teaching assistants, and students
should be collected and the materials corrected. Time required: 1
hour.
\end{enumerate}

{\bf Remark:}
Grade the good problem either as a significant portion of the homework or as a
quiz. A balance must be struck so the students will put effort into the
good problem without neglecting the others. 





